% main.tex — entry point, generated by the cours-latex skill.
% Compile with: ./compile.sh (Bash/WSL) or .\compile.ps1 (PowerShell).
% Placeholders <<MATIERE>>, <<SEMESTRE>>, <<AUTEUR>>, <<ETABLISSEMENT>> are
% filled in by the skill when it scaffolds a new course repo — angle
% brackets (not curly braces) on purpose, so the marker never collides
% with LaTeX's own {...} grouping syntax.
% oneside on purpose: this document is read on screen, not printed and
% bound — twoside mirrors page margins and (with the header setup in
% cours-preamble.sty) running-header sides for facing pages, which reads
% as a layout bug rather than a book-binding convention on screen.
\documentclass[french,11pt,oneside]{scrbook}
\usepackage{cours-preamble}

\title{<<MATIERE>>}
\author{<<AUTEUR>>}
\date{<<SEMESTRE>>}
\hypersetup{pdftitle={<<MATIERE>>}, pdfauthor={<<AUTEUR>>}}

\input{glossaire}

\begin{document}

\begin{titlepage}
  \centering
  \vspace*{2cm}
  {\sffamily\Huge\bfseries\color{coursnavy} <<MATIERE>> \par}
  \vspace{0.5cm}
  {\sffamily\Large\color{coursgray} <<ETABLISSEMENT>> \par}
  \vspace{2cm}
  {\Large <<SEMESTRE>> \par}
  \vfill
  {\large <<AUTEUR>> \par}
\end{titlepage}

\tableofcontents

% One \input per validated chapter — added by the skill as chapters are
% written, in source order.
% \input{chapitres/chap01-nom}

\printglossary[title=Glossaire, toctitle=Glossaire]

\end{document}
